\documentclass[11pt]{article}

\usepackage[T1]{fontenc}
\usepackage{lmodern}
\usepackage[margin=1in]{geometry}
\usepackage{amsmath,amssymb,amsthm}
\usepackage{microtype}
\usepackage[hidelinks]{hyperref}
\usepackage{url}

\newtheorem{theorem}{Theorem}
\newtheorem{lemma}[theorem]{Lemma}
\newtheorem{proposition}[theorem]{Proposition}
\newtheorem{corollary}[theorem]{Corollary}
\newcommand{\R}{\mathbb{R}}
\newcommand{\E}{\mathbb{E}}
\newcommand{\phat}{\widehat p}

\title{Log-convexity of the exercise boundary of an American put option with positive dividends}
\author{Robert Martin\thanks{The proof, formalisation, and manuscript were developed largely using
frontier AI systems (OpenAI GPT-6 Pro and Anthropic Fable 5.1). See \S\ref{sec:verification}
  and \hyperref[app:ai]{the appendix} for the verification scope and the use of AI.}\\
  Dianoa Labs}
\date{September 2026 --- \textbf{Draft}, not yet submitted or peer reviewed}

\begin{document}
\maketitle

\begin{abstract}
In the Black--Scholes model with constant risk-free interest rate $r>0$,
constant dividend yield $q$, and constant volatility $\sigma>0$, we prove
that the logarithm of the exercise boundary of an American put option is
convex in time-to-expiry whenever $0\le q\le r$. The proof
constructs an exact pricing-equation solution that matches the payoff
and smooth fit along a decreasing straight line. A maximum-principle
argument preserves the interval shape of the positive set of the option
price minus this comparison. An explicit backward-rectangle barrier then rules out
the contact configuration that would violate convexity. The geometric
argument does not require differentiability of the exercise boundary.
Established positive-time regularity gives nonnegative second derivatives
of the logarithmic boundary and strictly positive second derivatives
of the stock-price boundary as a separate corollary.
\end{abstract}

\section*{Introduction}

The exercise boundary of an American put option is the stock-price
threshold below which immediate exercise is optimal. We study how this
threshold varies with time-to-expiry in the Black--Scholes model. Although no general closed-form
expression for the boundary is known, its geometric properties depend
strongly on the dividend yield $q$ and the risk-free interest rate $r$.

For $q=0$, Ekstr\"om \cite{Ekstrom2004} and Chen, Chadam, Jiang,
and Zheng \cite{CCJZ2008} proved that the exercise boundary is a convex
function of time-to-expiry; the latter also
established strictly positive logarithmic curvature. For $q>r$,
Chen, Cheng, and Chadam \cite[Theorem~2]{CCC2013} showed that both
the boundary and its logarithm can fail to be convex, proving this
when $\log(q/r)$ is positive and sufficiently small. For positive
dividends below the risk-free interest rate, Liu's preprint \cite{Liu2017} presents
a proof of convexity of the stock-price boundary under
$r\ge q+\sigma^2/2$, where $\sigma$ is volatility, leaving
$0<q<r<q+\sigma^2/2$ unresolved.

We prove log-convexity throughout $0\le q\le r$, with $r>0$, covering
that remaining region and the endpoint $q=r$. The proof constructs an
exact pricing-equation solution fitted to a decreasing straight line
and uses maximum-principle arguments to exclude a contact configuration
incompatible with convexity. The geometric argument does not differentiate
the boundary; classical second-derivative conclusions follow separately
from established regularity.

\section{Main theorem and notation}

We work in the Black--Scholes model with constant parameters $r$, $q$,
and $\sigma$. Under the risk-neutral measure, the stock-price process satisfies
\[
  dS=(r-q)S\,du+\sigma S\,dW,\qquad S_0=S,
\]
where $u$ is elapsed time, $W$ is standard Brownian motion, and process
time subscripts are suppressed in the stochastic differential equation.
For strike $K>0$ and time-to-expiry $\tau>0$, the American put
option value is
\[
  P(S,\tau)=\sup_{0\le\theta\le\tau}
    \E_S\!\left[e^{-r\theta}(K-S_\theta)^+\right],
\]
where the supremum is over stopping times of the completed usual
Brownian filtration, with the horizon bound understood almost surely.
Write $B(\tau)$ for its exercise boundary, so that exercise is optimal
for $0<S\le B(\tau)$.

\begin{theorem}\label{thm:main}
Suppose $K>0$, $r>0$, $\sigma>0$, and $0\le q\le r$. Then
$\tau\mapsto\log(B(\tau)/K)$ is convex on $(0,\infty)$.
Moreover, $B$ is strictly convex on $(0,\infty)$.
\end{theorem}

Introduce normalized time-to-expiry $t$ and log-price $x$ by
\[
  a=\frac{\sigma^2}{2},\qquad t=a\tau,\qquad
  x=\log(S/K),\qquad k=\frac r a,\qquad h=\frac q a,\qquad
  \alpha=k-h-1.
\]
Define the normalized price $p(x,t)=P(Ke^x,t/a)/K$ and the
logarithmic exercise boundary $b(t)=\log(B(t/a)/K)$.
Thus $k>0$ and $0\le h\le k$. In the continuation region,
\begin{equation}\label{eq:pricing}
  p_t=p_{xx}+\alpha p_x-kp,\qquad x>b(t),\quad t>0,
\end{equation}
with value matching and continuation-side smooth fit
\begin{equation}\label{eq:fit}
  p(b(t),t)=1-e^{b(t)},\qquad p_x(b(t)+,t)=-e^{b(t)}.
\end{equation}
Here $p_x(b(t)+,t)$ denotes the right derivative of $x\mapsto p(x,t)$
at $b(t)$, not a derivative of the boundary.

We use standard properties independent of boundary convexity.
The price $p$ is continuous on $\R\times[0,\infty)$, is $C^{2,1}$
in continuation, and satisfies \eqref{eq:pricing} there, with
$p(x,0)=(1-e^x)^+$. The boundary $b$ is continuous on
$[0,\infty)$, with $b(0)=0$ and $b(t)<0$ for $t>0$.
At positive times, the price equals the payoff exactly when
$x\le b(t)$ and strictly exceeds it otherwise; value matching
and right spatial smooth fit hold as in \eqref{eq:fit}.
For $q=0$, see Chen--Chadam
\cite[\S2, especially Theorem~2.3 and Remark~2.2]{CC2007}; for
$q>0$, see Chen--Cheng--Chadam \cite[Theorem~1 and \S3]{CCC2013}.

Immediate exercise and exercise at expiry give
\[
  (1-e^x)^+\le p(x,t)\le1,\qquad p_E(x,t)\le p(x,t),
\]
where $p_E$ is the normalized European put value. A longer
time-to-expiry permits every stopping strategy available at a shorter
horizon, so $p$ is nondecreasing in $t$. Consequently, the exercise
region shrinks and $b$ is nonincreasing.

One elementary expiry estimate will be useful:
\begin{equation}\label{eq:expiry}
  \frac{b(t)}t\longrightarrow-\infty\qquad(t\downarrow0).
\end{equation}
Indeed, for every fixed $M>0$, the European put formula gives
$p_E(-Mt,t)=\sqrt{t/\pi}+O(t)$, whereas
$1-e^{-Mt}=O(t)$. Since the American value dominates the European
value, $-Mt$ lies in continuation for all sufficiently small $t$.
Hence $b(t)<-Mt$, which proves \eqref{eq:expiry}.

\section{An exact straight-line comparison}\label{sec:comparison}

Fix $c>0$ and $d<0$, and let $\ell(t)=d-ct$.
The construction treats the constant and exponential components of the
payoff separately. Both are eigenfunctions of the pricing operator
$\mathcal L=\partial_{xx}+\alpha\partial_x-k$, with
$\mathcal L1=-k$ and $\mathcal Le^x=-he^x$.
We therefore seek a comparison of the form
$f(x-\ell(t))-e^xg(x-\ell(t))$. Requiring it to solve the pricing
equation gives two constant-coefficient ordinary differential equations;
the initial conditions impose value matching and smooth fit along the
prescribed line. Define $f$ and $g$ by
\begin{equation}\label{eq:profiles}
  \begin{aligned}
    f''+(\alpha-c)f'-kf&=0,& f(0)&=1,& f'(0)&=0,\\
    g''+(\alpha+2-c)g'-hg&=0,& g(0)&=1,& g'(0)&=0.
  \end{aligned}
\end{equation}
For $\rho>0$, the solution of
$F''+\gamma F'-\rho F=0$, $F(0)=1$, $F'(0)=0$ is
\[
  F(z)=\frac{\lambda_+e^{\lambda_-z}-\lambda_-e^{\lambda_+z}}
              {\lambda_+-\lambda_-},\qquad
  \lambda_\pm=\frac{-\gamma\pm\sqrt{\gamma^2+4\rho}}2.
\]
Both coefficients are positive, so $F$ is convex and has its minimum
at zero. For $\rho=0$, the required solution is $F\equiv1$.
Consequently $f,g\ge1$ on $\R$.

With $z=x-\ell(t)$, define
\begin{equation}\label{eq:comparison}
  \phat(x,t)=f(z)-e^xg(z).
\end{equation}
Substitution into \eqref{eq:profiles} shows that $\phat$ solves
\eqref{eq:pricing} everywhere and satisfies
\[
  \phat(\ell(t),t)=1-e^{\ell(t)},\qquad
  \phat_x(\ell(t),t)=-e^{\ell(t)}.
\]
The comparison function need not be an option price.

\begin{lemma}\label{lem:domination}
For $t\ge0$ and $x\le0$, the comparison satisfies
$\phat(x,t)\ge1-e^x$.
\end{lemma}

\begin{proof}
Fix $t$, abbreviate $\ell=\ell(t)$, and let
\[
  H(z)=f(z)-1-e^{\ell+z}(g(z)-1).
\]
The initial values are $H(0)=H'(0)=0$, and direct calculation gives
\begin{equation}\label{eq:excess}
  H''+(\alpha-c)H'-kH
    =k-he^{\ell+z}+ce^{\ell+z}(g(z)-1).
\end{equation}
The right-hand side, denoted by $\Psi(z)$, is nonnegative whenever
$\ell+z\le0$, since $k\ge h\ge0$ and $g\ge1$.
Let $\lambda_-<0<\lambda_+$ be the roots for $f$. The Cauchy formula is
\[
  H(z)=\int_0^z E(z-y)\Psi(y)\,dy,\qquad
  E(s)=\frac{e^{\lambda_+s}-e^{\lambda_-s}}{\lambda_+-\lambda_-}.
\]
The kernel $E(s)$ has the sign of $s$. For $z\ge0$ the integral is
nonnegative. For $z<0$, the negative kernel and reversed integration
direction cancel. In both cases the integration interval stays below
the strike when $\ell+z\le0$, because $\ell<0$.
\end{proof}

Normalize the difference by the positive solution $f(z)$:
\[
  v(x,t)=\frac{p(x,t)-\phat(x,t)}{f(x-\ell(t))}.
\]
The quotient equation has no zero-order term:
\begin{equation}\label{eq:quotient}
  v_t=v_{xx}+\left(\alpha+2\frac{f'(z)}{f(z)}\right)v_x,
  \qquad v(b(t),t)\le0.
\end{equation}
The drift is bounded: $f'/f$ is a weighted average of the two
characteristic roots of $f$.

The far-field behavior is
\begin{equation}\label{eq:tail}
  v(x,t)\longrightarrow-1\qquad(x\to\infty),
\end{equation}
uniformly on bounded time intervals. To see this when $h>0$, let
$\lambda$ and $\mu$ be the positive roots for $f$ and $g$.
The polynomial $Q_f(s)=s^2+(\alpha-c)s-k$ satisfies
$Q_f(\mu+1)=-c<0$, so $\lambda>\mu+1$.
Thus $e^xg(z)/f(z)\to0$; also $p/f\to0$ since $p\le1$.
When $h=0$, use $g\equiv1$ and $Q_f(1)=-c$ instead.
The exponential estimates are uniform on each bounded time interval,
since $z=x-d+ct$ differs from $x$ by a uniformly bounded shift.

\section{Preservation of a positive interval}

The following three-point inequality preserves the interval shape of
superlevel sets at nonnegative levels.
Removing the zero-order term is key to this argument: at a maximizing
violating triple, the outer points are spatial maxima and the middle
point is a spatial minimum. The drift terms vanish, so the equation
makes the outer time derivatives nonpositive and the middle one
nonnegative. A small strict time penalty then makes these signs
incompatible with maximality.

\begin{lemma}[Parabolic three-point inequality]\label{lem:threepoint}
Let $T>0$ and let $\beta:[0,T]\to\R$ be continuous, with $\beta(t)<X$. Let
$u$ be continuous on the compact moving strip
\[
  Q=\{(x,t):0\le t\le T,\ \beta(t)\le x\le X\}.
\]
Suppose $u$ is $C^{2,1}$ at spatial interior points for $0<t\le T$
and solves $u_t=u_{xx}+A(x,t)u_x$ for some real-valued function $A$;
time derivatives at $T$ may be taken from the left. Assume
$u(\beta(t),t)\le0$ and $u(X,t)\le0$. If
\[
  \min\{u(x_1,0),u(x_3,0)\}\le\max\{u(x_2,0),0\}
  \qquad(\beta(0)\le x_1\le x_2\le x_3\le X),
\]
then, for every $0\le t\le T$,
\begin{equation}\label{eq:threepoint}
  \min\{u(x_1,t),u(x_3,t)\}\le\max\{u(x_2,t),0\}
  \qquad(\beta(t)\le x_1\le x_2\le x_3\le X).
\end{equation}
\end{lemma}

\begin{proof}
For $\delta>0$, define
\[
  P_\delta(s)=\frac{s+\sqrt{s^2+\delta^2}}2,\qquad
  M_\delta(s_1,s_2)=s_1-P_\delta(s_1-s_2).
\]
Their approximation bounds are
\[
  s^+\le P_\delta(s)\le s^++\delta/2,
  \qquad
  \min\{s_1,s_2\}-\delta/2\le M_\delta(s_1,s_2)\le\min\{s_1,s_2\}.
\]
Moreover, $0<P_\delta'<1$, so both partial derivatives of
$M_\delta$ are strictly positive.

For $\eta>0$, maximize the continuous function
\[
  G(x_1,x_2,x_3,t)
    =M_\delta(u(x_1,t),u(x_3,t))-P_\delta(u(x_2,t))-\eta t
\]
over the compact set $0\le t\le T$,
$\beta(t)\le x_1\le x_2\le x_3\le X$.
The bounds above give
\[
  G\le\min\{u(x_1,t),u(x_3,t)\}-\max\{u(x_2,t),0\}-\eta t.
\]
Thus $G\le0$ at $t=0$. At a collision $x_1=x_2$ or $x_2=x_3$,
the minimum is at most $u(x_2,t)$, so again $G\le0$.
If $x_1=\beta(t)$ or $x_3=X$, the minimum is nonpositive by the
lateral data. All other lateral contacts force a collision.
A positive maximum must therefore have
$t>0$ and $\beta(t)<x_1<x_2<x_3<X$.

Varying the spatial coordinates separately and using strict
monotonicity shows that $u(\cdot,t)$ has local maxima at $x_1,x_3$
and a local minimum at $x_2$. The first spatial derivatives vanish
there, so the drift terms disappear regardless of $A$, and the PDE gives
\[
  u_t(x_1,t)\le0,\qquad u_t(x_3,t)\le0,\qquad u_t(x_2,t)\ge0.
\]
Put $\rho=P_\delta'(u(x_1,t)-u(x_3,t))$ and
$\kappa=P_\delta'(u(x_2,t))$, both in $(0,1)$. Then
\[
  G_t=(1-\rho)u_t(x_1,t)+\rho u_t(x_3,t)-\kappa u_t(x_2,t)-\eta
      \le-\eta<0.
\]
Continuity of $\beta$ keeps this fixed spatial triple admissible at
all sufficiently close earlier times. Maximality therefore gives
$G_t\ge0$, including when $t=T$, a contradiction.
Thus $G\le0$; letting $\delta,\eta\downarrow0$ proves \eqref{eq:threepoint}.
\end{proof}

\begin{proposition}\label{prop:interval}
For every $t>0$ and $\varepsilon\ge0$, the set
$\{x>b(t):v(x,t)>\varepsilon\}$ is an interval or empty.
In particular, this holds for the positive set of $v$.
\end{proposition}

\begin{proof}
First check the initial shape. For $x>0$, we have $p(x,0)=0$.
Writing $z=x-d>0$ and
\[
  R(z)=e^{-z}\frac{f(z)}{g(z)},\qquad
  m=\frac{f'}f,\qquad n=\frac{g'}g,\qquad J=m-n-1,
\]
gives $v(x,0)=-1+e^d/R(z)$ and $R'=RJ$.
The profile equations give
\[
  m'=k-(\alpha-c)m-m^2,\qquad
  n'=h-(\alpha+2-c)n-n^2.
\]
At a zero of $J$, substitute $m=n+1$ and subtract. This gives
\begin{equation}\label{eq:crossing}
  J'=k-h-\alpha+c-1=c>0\qquad\text{at every zero of }J.
\end{equation}
Since $J(0)=-1$, every zero of $J$ is an upward crossing and there
is at most one: two upward crossings would require a transition
from positive to nonpositive values between them, contradicting
\eqref{eq:crossing}.
Thus $R$ is decreasing, or decreases and then increases, on $z>0$.
Its strict sublevel sets are intervals,
and hence every strict upper level set of $v(\cdot,0)$ is an interval.

This implies the initial three-point inequality: any violation
would place two endpoints strictly above a nonnegative level that
the middle value does not exceed. At $x=0$, it follows instead from
$v(0,0)\le0$, by payoff domination and the expiry payoff.

Fix $T>0$. By \eqref{eq:tail}, choose $X>0$ so large that
$v(X,t)<0$ for $0\le t\le T$. The other lateral value is nonpositive
by \eqref{eq:quotient}. Apply Lemma~\ref{lem:threepoint} directly to
$v$ on $b(t)\le x\le X$. Given any three ordered continuation
points, $X$ can also be chosen beyond the rightmost one. We obtain
\eqref{eq:threepoint} for $v$ at every positive time.

If the two endpoint values exceed $\varepsilon\ge0$, that inequality
forces the middle value to exceed $\varepsilon$ as well.
\end{proof}

\paragraph{Why the argument does not extend to larger dividends.}
For $q>r$, the limiting boundary at expiry is
$b(0+)=\log(r/q)<0$, and $b(t)<\log(r/q)$ for $t>0$
\cite[\S3.1 and Theorem~3]{CCC2013}. The comparison construction and
the identities \eqref{eq:quotient} and \eqref{eq:crossing} remain valid.
For lines with $d\le\log(r/q)$, the proof of
Lemma~\ref{lem:domination} still gives $\phat\ge1-e^x$ for
$x\le\log(r/q)$, since $k-he^x\ge0$ there. Thus the nonpositive
exercise-boundary data and the propagation argument survive.

The obstruction is the initial profile on the limiting continuation
domain $x>\log(r/q)$. On the additional segment $\log(r/q)<x\le0$,
$p(x,0)=1-e^x$, not zero, so
\[
  v(x,0)=\frac{1-e^x-\phat(x,0)}{f(x-d)}.
\]
Payoff domination is not established on this segment, and the
$R$-calculation controls only $x>0$. It therefore no longer proves
that the entire initial positive set is an interval. The propagation
lemma cannot rule out components already present initially. This
identifies the missing premise beyond $q=r$, consistently with the
Chen--Cheng--Chadam counterexamples; it is not an independent proof
of those counterexamples.

\section{Contact exclusion and geometric convexity}

We use the weak maximum and comparison principles and the strong
maximum principle in their classical parabolic forms; see
Lieberman \cite[Chapter~II]{Lieberman1996}.

\begin{lemma}[Backward-rectangle barrier]\label{lem:rectangle}
Let $L>0$ and $t_-<T$. Suppose $w$ is continuous on
$[0,L]\times[t_-,T]$, is $C^{2,1}$ on $(0,L)\times(t_-,T]$
(with left time derivatives at $T$), and satisfies
$w_t=w_{yy}+D(y,t)w_y$ there, with $D$ bounded.
Assume $w>0$ on the entire bottom and right edges, $w\ge0$ on the left
edge, and $w(0,T)=0$. Then
\[
  \liminf_{y\downarrow0}\frac{w(y,T)}y>0.
\]
\end{lemma}

\begin{proof}
Choose $A>\sup|D|$ and set
\[
  \phi(y)=\frac{e^{Ay}-1}{e^{AL}-1},\qquad
  \phi''+D\phi'=\frac{Ae^{Ay}}{e^{AL}-1}(A+D)>0.
\]
Since $\phi(0)=0$, $\phi(L)=1$, and the bottom and right edges
have positive minima, some $\eta>0$ satisfies $w\ge\eta\phi$
on the parabolic boundary. The difference $w-\eta\phi$ is a strict
supersolution, so weak comparison gives $w(y,T)\ge\eta\phi(y)$.
Consequently the displayed lower limit is at least
$\eta A/(e^{AL}-1)>0$.
\end{proof}


\begin{lemma}\label{lem:timeinterval}
For every $c>0$ and $d<0$, the set
$\{t>0:b(t)<d-ct\}$ is an interval or empty.
\end{lemma}

\begin{proof}
Suppose the assertion fails. There are times $t_1<t_m<t_2$ with the
boundary strictly below $\ell$ at $t_1$ and $t_2$, but not at $t_m$.
Continuity gives a first contact $t_0\in(t_1,t_m]$ such that
\[
  b(t)<\ell(t)\quad(t_1\le t<t_0),\qquad b(t_0)=\ell(t_0).
\]
There must be a point $x_*>b(t_0)$ with $v(x_*,t_0)>0$.
Otherwise $v(\cdot,t_0)\le0$ in continuation, and the weak maximum
principle would give $v\le0$ up to time $t_2$. Explicitly, choose
$X>\ell(t_2)$ large enough that $v(X,t)<0$ on $[t_0,t_2]$, using
\eqref{eq:tail}, and apply the principle on the bounded moving
strip $b(t)<x<X$ with its nonpositive initial and lateral data.
This contradicts $v(\ell(t_2),t_2)>0$,
since the line lies in continuation at $t_2$ and $\phat$ equals
the payoff there. The later return below the line is essential to
obtain this positive point on the contact slice.

Move with the line and set
\[
  w(y,t)=v(\ell(t)+y,t),\qquad L=x_*-\ell(t_0)>0.
\]
In a backward rectangle to the right of the line,
\begin{equation}\label{eq:moving}
  w_t=w_{yy}+D(y)w_y,\qquad
  D(y)=\alpha-c+2\frac{f'(y)}{f(y)}.
\end{equation}
Continuity gives $w(L,t)>0$ for $t$ sufficiently close to $t_0$.
For $t<t_0$ near $t_0$, we also have $w(0,t)>0$, since
$\ell(t)$ lies in continuation. Proposition~\ref{prop:interval}
then gives $w(y,t)>0$ for $0\le y\le L$.
Choose such a time $t_-<t_0$. On
$[0,L]\times[t_-,t_0]$, the bottom and right edges have strictly
positive data, and the left edge has nonnegative data. At contact,
value matching and smooth fit give
\begin{equation}\label{eq:contactfit}
  w(0,t_0)=0,\qquad w_y(0+,t_0)=0.
\end{equation}
Indeed, both $p-\phat$ and its right spatial derivative vanish at
contact, so the right derivative of their quotient by $f$ vanishes.

Lemma~\ref{lem:rectangle} now contradicts \eqref{eq:contactfit}.
\end{proof}

\begin{lemma}[No flat tail]\label{lem:noflattail}
The boundary cannot be constant on $[t_a,\infty)$ for any $t_a>0$.
\end{lemma}

\begin{proof}
Fix $\delta>0$. Let $f_0$ be the profile $f$ from
\eqref{eq:profiles} with $c=0$, so
$f_0''+\alpha f_0'-kf_0=0$, $f_0(0)=1$, $f_0'(0)=0$.
The profile formula in \S\ref{sec:comparison} gives $f_0\ge1$ and
bounded $f_0'/f_0$. Monotonicity of the price gives the nonnegative time increment
\[
  U(x,t)=\frac{p(x,t+\delta)-p(x,t)}{f_0(x)}\ge0.
\]
Nonincrease of $b$ ensures that both prices solve the pricing equation
when $x>b(t)$, so
\[
  U_t=U_{xx}+\left(\alpha+2\frac{f_0'}{f_0}\right)U_x.
\]
For $x>0$, $U(x,0)>0$: the expiry payoff is zero, whereas at time
$\delta$ the price strictly exceeds it. The strong maximum principle,
applied on finite cylinders in $x>0$, gives $U(x,t)>0$ for $x>0$,
$t>0$.

Suppose now that $b(t)=\bar b$ for $t\ge t_a>0$. Positivity at any
point above the strike at time $t_a$, followed by the same strong
maximum principle in $x>\bar b$, gives
$U(x,t)>0$ for $x>\bar b$ and $t>t_a$.
Value matching and smooth fit \eqref{eq:fit} at the common boundary give
$U(\bar b,t)=U_x(\bar b+,t)=0$ for $t\ge t_a$.

Choose $t_a<t_-<T$ and $L>0$, and set
\[
  w(y,t)=U(\bar b+T-t+y,t),\qquad
  0\le y\le L,\quad t_-\le t\le T.
\]
Its bottom and right edges are strictly positive, and its left edge
is nonnegative. It satisfies $w_t=w_{yy}+D(y,t)w_y$, where
\[
  D(y,t)=\alpha-1+
    2\frac{f_0'(\bar b+T-t+y)}{f_0(\bar b+T-t+y)}
\]
is bounded. Since $w(0,T)=w_y(0+,T)=0$,
Lemma~\ref{lem:rectangle} gives a contradiction.
\end{proof}

\begin{proof}[Proof of Theorem~\ref{thm:main}]
We first show that $b(t)/t$ is nondecreasing. If, for some $0<t_1<t_2$,
$b(t_2)/t_2<b(t_1)/t_1$, let $c=-b(t_1)/t_1>0$.
By \eqref{eq:expiry}, choose $t_*\in(0,t_1)$ with
$b(t_*)+ct_*<0$. Also $b(t_2)+ct_2<0$.
Choose
\[
  \max\{b(t_*)+ct_*,\ b(t_2)+ct_2\}<d<0.
\]
The boundary lies below $d-ct$ at $t_*$ and $t_2$, but above it at
$t_1$, where $b(t_1)+ct_1=0$. This contradicts
Lemma~\ref{lem:timeinterval}.

Now take any $0<t_1<t_2$ and write the chord through $(t_1,b(t_1))$ and
$(t_2,b(t_2))$ as $d-ct$. Nonincrease gives $c\ge0$.
If $c=0$, monotonicity makes $b$ constant on $[t_1,t_2]$, so the
chord inequality holds with equality. For $c>0$, the ratio
inequality just proved gives
\[
  d=\frac{t_2 b(t_1)-t_1 b(t_2)}{t_2-t_1}\le0.
\]
If $d<0$, raise the chord by any sufficiently small
$\varepsilon>0$ with $d+\varepsilon<0$.
Lemma~\ref{lem:timeinterval} puts the entire boundary segment
strictly below this raised chord. Letting $\varepsilon\downarrow0$
gives the chord inequality. If $d=0$, the same inequality follows
directly from $b(t)/t\le b(t_2)/t_2=-c$ for $t_1<t<t_2$.
Thus $b$ is convex, and positive rescaling of time proves the
logarithmic assertion.

We now deduce strict decrease. If $b(t_1)=b(t_2)$ for some $0<t_1<t_2$,
convexity gives, for every $t>t_2$,
\[
  0=\frac{b(t_2)-b(t_1)}{t_2-t_1}\le\frac{b(t)-b(t_2)}{t-t_2}\le0,
\]
where the last inequality is monotonicity. Thus $b$ has a flat
tail, contrary to Lemma~\ref{lem:noflattail}.
Consequently $b(t_1)\ne b(t_2)$ whenever $t_1\ne t_2$.
For $0<\theta<1$, convexity of $b$ and strict convexity of the
exponential give
\[
  e^{b((1-\theta)t_1+\theta t_2)}
    \le e^{(1-\theta)b(t_1)+\theta b(t_2)}
    <(1-\theta)e^{b(t_1)}+\theta e^{b(t_2)}.
\]
Since $B(\tau)=Ke^{b(a\tau)}$, this proves strict convexity of $B$.
\end{proof}

Since $b$ is continuous at $0$ with $b(0)=0$, its convexity extends
to $[0,\infty)$. By \eqref{eq:expiry}, its right derivative at $0$
is $-\infty$.

\section{Classical second derivatives}\label{sec:classical}

The positive-time regularity results of Chen--Chadam
\cite[Remark~2.2 and subsequent regularity analysis]{CC2007} for zero
dividends, and Chen--Cheng \cite{ChenCheng2012} for dividend-paying assets
(see also Chen--Cheng--Chadam \cite[Theorem~1 and \S3]{CCC2013}),
give $b\in C^2((0,\infty))$, in fact
$C^\infty$. They do not require boundary convexity.

\begin{corollary}
Under the hypotheses of Theorem~\ref{thm:main}, for every $\tau>0$,
\[
  \frac{d^2}{d\tau^2}\log\frac{B(\tau)}K\ge0,
  \qquad B''(\tau)>0.
\]
\end{corollary}

\begin{proof}
The symmetric second differences of a convex $C^2$ function are
nonnegative, so $b''\ge0$. Convexity and strict decrease also give
\[
  b'(t)\le\frac{b(t+s)-b(t)}s<0\qquad(s>0).
\]
Twice differentiating $B(\tau)=Ke^{b(a\tau)}$ yields
\[
  B''(\tau)=Ka^2e^{b(a\tau)}
       \bigl(b''(a\tau)+b'(a\tau)^2\bigr)>0.
\]
The logarithmic second derivative is $a^2b''(a\tau)\ge0$.
\end{proof}

\section{Formal verification}\label{sec:verification}

The accompanying Lean development \cite{Formalization} proves the
geometric theorem and, separately, $C^2$ regularity and the classical
curvature corollary. Its price is defined by optimal stopping over
all bounded stopping times of the completed usual Brownian filtration;
a proved equivalence permits almost-sure horizon bounds. The normalized
price uses strike one and volatility $\sqrt2$, with the boundary defined
by payoff contact. The pricing equation, smooth fit, and identification
with the physical boundary are proved from these definitions, including
$0<B(\tau)<K$ for $\tau>0$.

The paper follows Lean's comparison, three-point maximum, rectangle
barrier, and geometric arguments. Both use nonincrease to prove
convexity, then exclude a flat tail to obtain strict decrease.
Two analytic steps are presented differently: the paper derives
\eqref{eq:expiry} from the European put asymptotic, while Lean uses
a local square-root subsolution; the positivity propagation in
Lemma~\ref{lem:noflattail} is proved in Lean by explicit local barriers.

The development builds on Mathlib \cite{Mathlib2020}, the BrownianMotion library
of Degenne, Ledvinka, Marion, and Pfaffelhuber \cite{Degenne2025},
and Coelho's MathFin library \cite{Coelho2026}. MathFin supplies
geometric Brownian motion, It\^o and local-martingale results,
heat-kernel tools, and filtration machinery. The boundary-convexity
argument is proved in this project, not imported from these libraries.

The recorded dependency audit rechecks the full proof closure in a
fresh Lean kernel. It reports only the standard axioms
\texttt{propext}, \texttt{Classical.choice}, and \texttt{Quot.sound},
with no unproved project axioms or \texttt{sorry} dependencies.
Kernel checking certifies the formal statements and proofs; their
correspondence with the intended financial model and this manuscript
remains a separate mathematical check.

\appendix
\section*{Appendix: AI usage}\label{app:ai}

The proof development, formalization, and manuscript preparation were
carried out largely by frontier AI models, with the author selecting the problem, directing the investigation, and making editorial decisions.

While researching American options, the author encountered the
unresolved boundary-convexity question for the case of nonzero dividends and considered
it a promising subject for AI-assisted mathematical discovery, because the problemw as relatively simple to state
and a counterexample would be easily verifiable.

The first task given to ChatGPT-6 Pro (Astra) was to find such a
counterexample, and to provide motivation to the model, the author suggested to the model that he had already found one.\footnote{For a related discussion of motivational prompting, see Anthropic,
\href{https://www.anthropic.com/research/riemann-zeta}{``Learning more about Claude's mathematical capabilities''},
August 10, 2026, especially ``Claude's methodology.''}

After an extended search, Astra had not found one
and instead proposed structural reasons to expect convexity.
The author then suggested that his proof of a counterexample may have been erroneous, and asked it to instead prove the full conjecture in the
nonnegative-dividend regime $q\le r$.

To the Author's surpirse, Astra produced a candidate proof in 31 minutes. The author asked Anthropic's
Fable 5.1 to examine it for errors, and Fable reported none.
The project then moved to the Codex command-line environment for
formal verification in Lean, performed primarily by GPT-6 Astra (Pro),
with Fable 5.1 (High) in Claude Code providing additional review during the development.

Extensive prompting was used to build confidence that the proof and its formalization were correct, including prompts
to deeply analyze the extant literature to understand why this approach had not been discovered before,
and explicit prompting to consider and avoid the various ways a Lean formalization could be falsely constructed.

The manuscript was drafted by Codex (prompted to follow good technical writing practices\footnote{Stephen Boyd, Ernest K. Ryu, and Neal Parikh,
\href{https://web.stanford.edu/class/ee364b/latex_templates/template_notes.pdf}{``\LaTeX{} Style Guide for EE 364B''},
May 12, 2014.}), reviewed by Fable 5.1, then revised and edited manually by the author.


\bibliographystyle{plain}
\begingroup
\small
\raggedright
\bibliography{references}
\endgroup

\end{document}
